% !TEX root = ./main.tex
\chapter{Architettura e Tecnologie Utilizzate}
\label{chap:architettura}

\section{Requisiti del sistema}
\label{sec:requisiti}

Un prototipo di agente conversazionale embodied come SOULFRAME mette insieme interazione in tempo reale, gestione di asset 3D e servizi di inferenza esterni. In questa sezione i requisiti servono soprattutto a chiarire quali capacità erano davvero necessarie per far funzionare il prototipo e quali vincoli hanno inciso sulla sua tenuta tecnica.

\subsection{Requisiti funzionali}

Le funzioni principali del sistema sono distribuite tra servizi FastAPI \cite{FastAPI2026} separati, richiamati dal client Unity WebGL tramite HTTP. L'elenco seguente mantiene il formato dei requisiti funzionali, ma va letto come descrizione operativa del prototipo più che come capitolato formale.

RF1 (Libreria e selezione avatar): il client deve ottenere e visualizzare la lista degli avatar disponibili tramite \texttt{/avatars/list} e permettere la selezione di un profilo attivo identificato da \texttt{avatar\_id}. Il requisito è soddisfatto se la lista include almeno un modello locale di fallback definito in \texttt{LOCAL\_MODELS} e, quando il backend è raggiungibile, include anche avatar importati con un URL caricabile del modello \texttt{.glb}.

RF2 (Import e caching asset \texttt{.glb}): quando è disponibile un URL di export (ad es. da Avaturn \cite{AvaturnDocs2026}), il client deve richiedere l'import con \texttt{/avatars/import} e ottenere un URL di cache server-side da usare per il download e l'istanza in scena. Il requisito è soddisfatto se import ripetuti dello stesso URL non creano duplicati, mantengono lo stesso \texttt{avatar\_id} con URL di cache stabile, e se l'avatar resta caricabile anche dopo riavvio dei servizi.

RF3 (Setup voce persistente per avatar): il client deve consentire la registrazione di un campione vocale e inviarlo al TTS con \texttt{/set\_avatar\_voice}, associandolo all'avatar attivo. Il requisito è soddisfatto se la registrazione supera una verifica di similarità testuale con soglia definita e se il riferimento vocale risulta riutilizzabile nelle sessioni successive, insieme alle frasi di attesa già preparate per l'avatar.

RF4 (Memoria per-avatar con ingest e retrieval contestuale): il sistema deve salvare note e contenuti da file nella memoria dell'avatar tramite \texttt{/remember} e \texttt{/ingest\_file}, e deve usare tale memoria nella generazione contestuale via \texttt{/chat}. Il requisito è soddisfatto se la memoria resta persistente tra sessioni, se due avatar distinti non condividono contenuti indicizzati e se la risposta del servizio espone metadati minimi di tracciabilità coerenti con il turno.

RF5 (Turno vocale end-to-end in push-to-talk): il client deve gestire una conversazione a turni acquisendo audio, inviandolo allo STT con \texttt{/transcribe}, inoltrando la trascrizione al RAG con \texttt{/chat} e riproducendo la risposta audio tramite streaming con \texttt{/tts\_stream}. Il requisito è soddisfatto se, per ogni turno, il sistema produce testo trascritto, testo di risposta e avvio del playback audio con transizioni UI coerenti tra ascolto, elaborazione e riproduzione.

RF6 (Ingest multimodale per memoria per-avatar): il sistema deve consentire l'ingestione multimodale nella memoria dell'avatar, includendo descrizione immagine con \texttt{/describe\_image} e ingestione documentale con \texttt{/ingest\_file} (estrazione e indicizzazione del contenuto). Il requisito è soddisfatto se, dopo l'operazione, i contenuti risultano recuperabili tramite \texttt{/recall} e/o se \texttt{/avatar\_stats} indica che per l'avatar è presente una memoria già popolata.

\subsection{Requisiti non funzionali}

I requisiti non funzionali sono formulati secondo il modello ISO/IEC 25010 e si concentrano sugli aspetti che, nel prototipo, incidono di più sulla plausibilità dell'interazione e sulla sua gestione tecnica. In VR anche modifiche piccole possono avere effetti evidenti su tempi, stabilità e leggibilità dell'esperienza; per questo è utile esplicitare fin da qui i vincoli di qualità che hanno guidato le scelte architetturali \cite{Karre2024VReqST}.

RNF1 (Efficienza prestazionale - time behaviour): la latenza percepita deve restare controllata misurando il tempo tra rilascio del comando push-to-talk e primo campione audio riprodotto dal TTS. Nel prototipo è stata adottata come soglia operativa una risposta udibile entro 12 s in esecuzione locale e entro 16 s in deploy pubblico. La verifica, in questa fase, resta osservativa: i timestamp lato client sono utili per diagnosi e confronto tra ambienti, mentre una telemetria turn-based persistente resta necessaria per una validazione quantitativa completa, come discusso nel Capitolo~\ref{chap:risultati}. Le soglie 12/16 s vanno quindi lette come target di progetto, non come benchmark conclusivo \cite{Lin2024Estuary}.

RNF2 (Affidabilità - availability e fault tolerance): durante sessioni prolungate il sistema deve mantenere utilizzabile la maggior parte dei turni vocali, dove un turno è riuscito se STT, \texttt{/chat} e \texttt{/tts\_stream} restituiscono esito valido entro timeout. Sono implementati timeout, ritentativi limitati e gestione esplicita degli errori lato client; in caso di fallimento il flusso degrada su messaggi di stato UI senza blocchi permanenti. Nel prototipo la verifica resta soprattutto osservativa, riletta attraverso log di sessione, errori HTTP e comportamento della macchina a stati.

RNF3 (Manutenibilità, modularità e portabilità): la sostituzione di un micro-servizio deve richiedere solo una variazione della configurazione centralizzata lato client, mantenendo invariati endpoint e dati attesi. Il requisito si considera soddisfatto quando il frontend WebGL continua a funzionare sia nello sviluppo locale, con URL assoluti, sia in deploy dietro reverse proxy con percorsi del tipo \texttt{/api/<servizio>}, e quando dopo un aggiornamento viene rieseguita una checklist minima di compatibilità.

RNF4 (Coerenza estetica e comunicazione di stato): il linguaggio visivo deve restare coerente con la natura stilizzata degli avatar Avaturn e con il carattere sperimentale del prototipo. Nel client questo passa da effetti di post-processing, rings animati e feedback di selezione, scelti perché compatibili con rendering WebGL e utili a rendere leggibili gli stati operativi senza aumentare il carico cognitivo dell'utente. Il requisito è soddisfatto se gli indicatori visivi accompagnano in modo consistente transizioni e operazioni principali dell'interfaccia (inizializzazione, setup, onboarding, selezione), mantenendo continuità tra resa dell'avatar e contesto UI.

\section{Architettura di riferimento di SOULFRAME}
\label{sec:architettura-riferimento}
\subsection{Vista d'insieme dei componenti frontend/backend}
SOULFRAME adotta un'architettura client--server a componenti, pensata per lasciare al browser ciò che dipende da interfaccia e resa 3D e demandare al backend inferenza e persistenza. Il client è una build Unity WebGL eseguita nel browser: gestisce la UI, le transizioni di stato dell'esperienza e l'interazione \textit{push-to-talk}, oltre al rendering e alla visualizzazione dell'avatar. La comunicazione verso il backend avviene via HTTP ed è mediata da una configurazione centralizzata, così che il client possa funzionare sia in esecuzione locale sia in deploy pubblico senza cambiare logica applicativa. L'intero impianto resta orientato a una conversazione individuale per sessione: memoria persistente e contesto recente vengono organizzati per \texttt{avatar\_id} e \texttt{session\_id}, ma non esiste una sincronizzazione reale tra utenti concorrenti.

Sul backend, la pipeline è scomposta in quattro micro-servizi FastAPI indipendenti. Whisper trasforma l'audio in testo, il servizio RAG governa orchestrazione conversazionale e memoria per-avatar, il TTS genera la voce dell'avatar e l'Avatar Asset Server mantiene il catalogo dei modelli \texttt{.glb}. I dettagli delle route compaiono più avanti e nelle appendici; nel corpo del capitolo interessa soprattutto mostrare come queste responsabilità restino ben distinte.

Le route principali riflettono bene questa divisione. Whisper espone \texttt{/transcribe}; il servizio RAG concentra \texttt{/chat}, \texttt{/chat\_session/start}, \texttt{/remember}, \texttt{/recall} e ingestione dei contenuti; il TTS offre sintesi completa, sintesi in streaming e gestione della voce per-avatar; il servizio asset governa \texttt{/avatars/list}, \texttt{/avatars/import} e il serving dei modelli \texttt{.glb}. Più che seguire ogni payload, qui è utile notare che ciascun blocco conserva un perimetro abbastanza chiaro.

Un runtime separato, Ollama, fornisce modelli di chat ed embeddings. In questo modo il servizio RAG resta focalizzato sull'orchestrazione applicativa e i modelli possono essere aggiornati o sostituiti senza toccare il client, purché restino stabili interfacce e formati attesi.

Nel prototipo questa impostazione si riflette anche nelle configurazioni dei due ambienti. In locale Windows Whisper lavora con il profilo \texttt{small}, mentre nell'ambiente Ubuntu usato per le prove è stato possibile adottare \texttt{medium}; sul lato chat, la variante usata via Ollama cambia tra sviluppo e deploy, ma questi dettagli restano confinati al backend.

In deploy pubblico, l'elemento chiave di integrazione è il reverse proxy Caddy, con due funzioni: terminazione TLS e punto d'ingresso unico per il browser. Il client WebGL comunica con un'unica origine HTTPS e invoca le API tramite percorsi riscritti (\texttt{/api/whisper/*}, \texttt{/api/rag/*}, \texttt{/api/avatar/*}, \texttt{/api/tts/*}). Caddy inoltra le richieste alle porte interne 8001--8004, mantenute non esposte verso l'esterno.

La soluzione riduce la complessità lato browser, isola i servizi e impedisce che una build pubblica tenti di contattare \texttt{127.0.0.1:800x} dal browser dell'utente finale. Nel prototipo, gli endpoint non implementano autenticazione o autorizzazione applicativa; l'architettura presuppone quindi una rete di esecuzione controllata.

\begin{figure}[H]
\centering
\includegraphics[width=1.0\textwidth]{pipeline_e2e.png}
\caption{Flusso end-to-end del sistema SOULFRAME: dall'input vocale dell'utente alla risposta dell'agente embodied.}
\label{fig:pipeline-e2e}
\end{figure}

Figura~\ref{fig:arch-cap3} riassume questa vista a blocchi. Figura~\ref{fig:pipeline-e2e} mostra in forma sintetica la catena conversazionale completa, dall'input vocale alla risposta dell'agente embodied. Le due figure sono complementari: la prima mette a fuoco il flusso logico della conversazione, la seconda come lo stesso flusso viene distribuito tra client, proxy e servizi.

La stessa scomposizione in componenti viene mantenuta sia in locale sia su server. In ambiente Windows lo script di avvio coordina l'esecuzione dei servizi e applica controlli di base, per esempio sulla disponibilità delle porte, così da ridurre conflitti tra processi durante lo sviluppo. In ambiente Ubuntu, ciascun micro-servizio è gestito come unità \texttt{systemd}, con comandi amministrativi che permettono start, stop, status e aggiornamenti senza riconfigurare manualmente l'intero stack. È una soluzione più onerosa di un eseguibile unico, ma nel prototipo ha reso più leggibili i problemi che emergevano tra servizi diversi.

\begin{figure}[H]
\centering
\includegraphics[width=1.0\textwidth]{arch_cap3.png}
\caption{Vista d'insieme dell'architettura a componenti di SOULFRAME: il client Unity WebGL comunica con i micro-servizi backend attraverso il reverse proxy Caddy.}
\label{fig:arch-cap3}
\end{figure}

\subsection{\texorpdfstring{Flusso end-to-end audio $\rightarrow$ testo $\rightarrow$ risposta $\rightarrow$ audio}{Flusso end-to-end audio -> testo -> risposta -> audio}}
Il flusso conversazionale end-to-end segue una sequenza a turni: parte dall'input vocale dell'utente e termina con la riproduzione della risposta sintetizzata, mantenendo sul client Unity WebGL le responsabilità di interazione e rendering e delegando al backend le fasi di inferenza. L'interazione è di tipo push-to-talk: l'utente tiene premuto \texttt{SPACE} per parlare e, al rilascio, il client chiude la registrazione e prepara un file audio WAV. In ambiente WebGL l'acquisizione del microfono richiede un bridge JavaScript, realizzato qui tramite \texttt{AudioCapture.jslib}.\footnote{Unity Technologies, \emph{Audio in WebGL}, \url{https://docs.unity3d.com/6000.2/Documentation/Manual/webgl-audio.html}}

Una volta completata la cattura, il client invia l'audio al servizio STT tramite \texttt{/transcribe}. Questa prima trasformazione da segnale audio a testo rende poi possibile l'intero passaggio conversazionale successivo.

Il servizio RAG/LLM riceve la trascrizione tramite \texttt{/chat}, includendo l'identificativo dell'avatar (\texttt{avatar\_id}). Qui il backend classifica la richiesta, sceglie il percorso di recupero più adatto e costruisce il contesto della risposta.

Prima del primo turno, però, il client può aprire anche una sessione esplicita con \texttt{/chat\_session/start}. Questo passaggio genera un \texttt{session\_id} e permette di tenere distinti memoria persistente dell'avatar, contesto recente e traccia della sessione.

La memoria per-avatar è costruita su un database vettoriale persistente (ChromaDB \cite{ChromaDB2026}), popolato tramite embedding generati con un modello dedicato. Accanto a questo livello persistente, il backend mantiene anche una finestra scorrevole dei turni più recenti in RAM, indicizzata per coppia \texttt{(avatar\_id, session\_id)}. Questa finestra non sostituisce la memoria interrogata via retrieval: conserva continuità locale tra richieste ravvicinate, mentre il log di sessione su file resta distinto. Il servizio invoca poi il modello linguistico per la risposta e, quando necessario, applica una riscrittura con vincoli di aderenza ai dati recuperati.

La scelta di tenere insieme questi tre livelli --- memoria persistente, contesto recente e log --- nasce da un compromesso preciso. Una memoria solo persistente rende più difficile seguire riferimenti ravvicinati; una memoria solo di sessione fa perdere continuità tra incontri diversi; un semplice log non basta a sostenere il turno in tempo reale. Nel corpo interessa soprattutto questa distinzione, mentre i meccanismi più minuti di retrieval e tracciabilità sono sviluppati nel Capitolo~\ref{chap:implementazione} e in Appendice~\ref{app:rag}.

In questa tesi tali scelte sono trattate come euristiche implementative, non come prova statistica definitiva dell'aderenza ai dati recuperati. È anche il punto della pipeline in cui si sommano più facilmente dipendenza dal contesto e variabilità computazionale; per questo la latenza percepita entra direttamente nel disegno architetturale.

Ricevuto il testo di risposta, il client attiva la sintesi vocale tramite il servizio Text-to-Speech (TTS) basato su Coqui XTTS v2. La richiesta avviene preferibilmente tramite l'endpoint di streaming \texttt{/tts\_stream}, includendo \texttt{avatar\_id} e lingua: il servizio recupera il profilo vocale associato all'avatar (voice cloning) e produce audio progressivamente, consentendo al client di iniziare il playback non appena arrivano i primi byte dello stream. In pratica, l'attesa viene spostata dalla generazione dell'intero file alla sola apertura della risposta udibile, con un effetto diretto sulla continuità percepita del turno.

Per mitigare i tempi morti più evidenti, il TTS usa una pre-attivazione iniziale e, durante i turni, ricorre a brevi \emph{wait phrases} prima dell'avvio della risposta principale in streaming. Figura~\ref{fig:ui-state-diagram-cap3} colloca questo comportamento nel ciclo degli stati UI; i dettagli sono approfonditi nel Capitolo~\ref{chap:implementazione}.

Conviene anche distinguere i due tipi di output vocale. Le \emph{wait phrases} sono clip brevi, già associate all'avatar dopo il setup voce e usate come segnale immediato di continuità; la risposta principale, invece, arriva come stream WAV progressivo, con un costo maggiore ma con contenuto semantico completo. Nel prototipo questa distinzione ha inciso sulla qualità percepita del turno più di quanto possa sembrare da una lettura solo tecnica.

\begin{figure}[H]
\centering
\includegraphics[width=0.96\textwidth]{ui_state_diagram_cap3.png}
\caption{Diagramma degli stati UI del client SOULFRAME: dalla fase
di caricamento iniziale al ciclo conversazionale
(Idle $\rightarrow$ Listening $\rightarrow$ Processing $\rightarrow$ Speaking).}
\label{fig:ui-state-diagram-cap3}
\end{figure}

Figura~\ref{fig:flusso-e2e-cap3} sintetizza la sequenza temporale e i formati dati scambiati tra componenti, evidenziando le trasformazioni principali (WAV$\rightarrow$testo, testo$\rightarrow$testo contestualizzato, testo$\rightarrow$audio).

Nel blocco RAG la figura distingue i tre passaggi essenziali: capire la richiesta, recuperare il contesto utile e controllare la risposta finale. In deploy con reverse proxy, questi passaggi restano invariati anche se le route vengono raggiunte tramite percorsi riscritti sotto un'unica origine HTTPS.

\begin{figure}[H]
\centering
\includegraphics[width=1.0\textwidth]{flusso_e2e_cap3.png}
\caption{Flusso end-to-end di una richiesta conversazionale in SOULFRAME: dall'acquisizione audio push-to-talk alla riproduzione della risposta vocale, con dettaglio dei passaggi interni RAG (capire la richiesta, recuperare contesto, controllare la risposta).}
\label{fig:flusso-e2e-cap3}
\end{figure}

\subsection{Flusso di gestione avatar (creazione, import, cache, rendering)}
Il flusso di gestione avatar affianca due percorsi che convergono nella stessa esperienza in scena: avatar locali pre-inclusi nella build e avatar importati dall'utente. Il punto non è solo offrire due modalità di accesso, ma garantire una libreria sempre utilizzabile e, insieme, un canale controllato per acquisire e servire asset esterni in formato \texttt{.glb}.

Nel percorso locale, il client seleziona un profilo già disponibile e lo carica direttamente in scena. Nel percorso importato, invece, il profilo viene creato nell'overlay web di Avaturn e il client delega al backend l'import del \texttt{.glb}, ottenendo un \texttt{cached\_glb\_url} da cui caricare il modello. Da quel punto in poi il flusso torna unico: selezione del profilo, caricamento del modello e attivazione dell'avatar nella conversazione.

Dopo selezione o creazione, il flusso prosegue con onboarding voce e memoria per-avatar. L'accesso a \texttt{MainMode} avviene solo quando i prerequisiti minimi del profilo risultano soddisfatti, così da evitare conversazioni prive di contesto operativo. La presenza di modelli locali di fallback evita anche che la libreria si svuoti quando il backend avatar non è raggiungibile. La Figura~\ref{fig:flusso-avatar-cap3} esplicita questo raccordo, mentre i dettagli di import e robustezza lato asset sono approfonditi nel Capitolo~\ref{chap:implementazione}.

\begin{figure}[!htbp]
\centering
\includegraphics[width=1.0\textwidth]{flusso_avatar_cap3.png}
\caption{Flusso di gestione avatar in SOULFRAME: dalla creazione tramite Avaturn alla cache server-side e al rendering nel client Unity.}
\label{fig:flusso-avatar-cap3}
\end{figure}

\begin{figure}[!htbp]
\centering
\includegraphics[width=0.95\textwidth]{ui_carosello.png}
\caption{Interfaccia carosello avatar in SOULFRAME: selezione del
profilo attivo tra avatar locali e importati.}
\label{fig:ui-carosello-cap3}
\end{figure}

\begin{figure}[!htbp]
\centering
\includegraphics[width=0.95\textwidth]{ui_mainmenu.png}
\caption{Interfaccia \texttt{MainMenu} in SOULFRAME: accesso alle funzioni principali prima dell'onboarding dell'avatar.}
\label{fig:ui-mainmenu-cap3}
\end{figure}

\begin{figure}[!htbp]
\centering
\includegraphics[width=0.95\textwidth]{ui_setupvoice.png}
\caption{Schermata \texttt{SetupVoice}: registrazione e validazione del campione vocale durante l'onboarding.}
\label{fig:ui-setupvoice-cap3}
\end{figure}

\begin{figure}[!htbp]
\centering
\includegraphics[width=0.95\textwidth]{ui_setupmemory.png}
\caption{Schermata \texttt{SetupMemory}: inserimento note e ingestione contenuti per popolare la memoria per-avatar.}
\label{fig:ui-setupmemory-cap3}
\end{figure}

\begin{figure}[!htbp]
\centering
\includegraphics[width=0.95\textwidth]{ui_push_to_talk.png}
\caption{Stato di ascolto push-to-talk nel client Unity WebGL:
acquisizione audio attiva con indicatore visivo.}
\label{fig:ui-ptt-cap3}
\end{figure}

Le Figure~\ref{fig:ui-carosello-cap3}, \ref{fig:ui-mainmenu-cap3}, \ref{fig:ui-setupvoice-cap3}, \ref{fig:ui-setupmemory-cap3} e \ref{fig:ui-ptt-cap3} rendono leggibile la sequenza concreta di accesso al sistema: scelta del profilo, ingresso nel menu, onboarding voce, onboarding memoria e avvio del turno conversazionale.

\FloatBarrier

\section{Componenti software implementati}
I componenti implementati rendono operativa l'architettura della Sezione~\ref{sec:architettura-riferimento}: integrazione di Avaturn nel client Unity WebGL, backend AI a micro-servizi e persistenza dei dati legati all'avatar. La panoramica unificata dei servizi backend è riportata in Tabella~\ref{tab:microservizi-cap3}; i dettagli implementativi più minuti sono nel Capitolo~4.

\subsection{Componenti principali del prototipo}
Il prototipo prende forma dall'unione tra integrazione di Avaturn nel client Unity WebGL e backend AI a micro-servizi. L'avatar viene creato in ambiente web dedicato, il client ne acquisisce l'identità, delega l'import del modello \texttt{.glb} al backend e poi lo riporta nel percorso unitario di SOULFRAME descritto in Sezione~\ref{sec:architettura-riferimento}. Nel runtime Unity il caricamento dei \texttt{.glb} avviene con GLTFast\footnote{Unity Technologies, \emph{glTFast package documentation}, \url{https://docs.unity3d.com/Packages/com.unity.cloud.gltfast@6.10/manual/index.html}}, così da tenere la creazione dell'avatar fuori dalla scena ma il suo uso operativo dentro lo stesso flusso applicativo.

Sul backend, invece, SOULFRAME è realizzato come insieme di micro-servizi indipendenti basati su FastAPI, avviati su porte dedicate. La scomposizione in servizi rende esplicita la separazione delle responsabilità e permette di verificare disponibilità e corretto instradamento con controlli semplici (\texttt{/health}) usati sia in fase operativa sia nell'avvio iniziale del client. Tabella~\ref{tab:microservizi-cap3} sintetizza i quattro componenti e i relativi punti di accesso.

{\scriptsize
\setlength{\tabcolsep}{2pt}
\renewcommand{\arraystretch}{1.2}
\begin{longtable}{|>{\raggedright\arraybackslash}p{0.16\textwidth}|>{\raggedright\arraybackslash}p{0.20\textwidth}|>{\raggedright\arraybackslash}p{0.07\textwidth}|>{\raggedright\arraybackslash}p{0.24\textwidth}|>{\raggedright\arraybackslash}p{0.27\textwidth}|}
\caption{Micro-servizi backend di SOULFRAME: porte, endpoint principali e responsabilità operative usate nel prototipo.}
\label{tab:microservizi-cap3}\\
\hline
\textbf{Componente} & \textbf{File Python} & \textbf{Porta} & \textbf{Endpoint chiave} & \textbf{Responsabilità} \\
\hline
\endfirsthead
\hline
\textbf{Componente} & \textbf{File Python} & \textbf{Porta} & \textbf{Endpoint chiave} & \textbf{Responsabilità} \\
\hline
\endhead
\hline
\endfoot
\hline
\endlastfoot
Whisper STT & \path{whisper_server.py} & \texttt{8001} & \path{GET /health}\ \path{POST /transcribe} & Trascrizione Speech-to-Text (STT) di audio caricato (WAV o formati compatibili), gestione lingua e cleanup di file temporanei. \\
\hline
RAG/LLM & \path{rag_server.py} & \texttt{8002} & \path{GET /health}\ \path{GET /avatar_stats}\ \path{POST /chat}\ \path{POST /chat_session/start}\ \path{POST /remember}\ \path{POST /recall}\ \path{POST /ingest_file}\ \path{POST /describe_image} & Costruzione della risposta contestuale: memoria per-avatar su ChromaDB, classificazione dell'intento, ricerca mista lessicale+semantica, filtraggio dei passaggi più affidabili, controllo finale del testo e metadati di turno (\texttt{intent}, \texttt{rag\_used}, \texttt{auto\_remembered}). \\
\hline
Coqui XTTS TTS & \path{coqui_tts_server.py} & \texttt{8004} & \path{GET /health}\ \path{POST /tts}\ \path{POST /tts_json}\ \path{POST /set_avatar_voice}\ \path{GET/DELETE /avatar_voice}\ \path{GET /wait_phrase} & Sintesi Text-to-Speech (TTS) con voice cloning per avatar, streaming audio e serving di frasi di attesa già predisposte nel setup voce; gestione e persistenza dei profili vocali per \texttt{avatar\_id}. \\
\hline
Avatar Asset Server & \path{avatar_asset_server.py} & \texttt{8003} & \path{GET /health}\ \path{GET /avatars/list}\ \path{POST /avatars/import}\ \path{GET /avatars/{id}/model.glb}\ \path{DELETE /avatars/{id}} & Import e caching server-side di modelli \texttt{.glb}, deduplicazione via hash URL, riallineamento automatico limitato dei metadati e lista unificata di avatar locali di fallback e avatar importati. \\
\end{longtable}
}

Nel flusso reale i due piani restano vicini ma non coincidono. Quando l'avatar viene creato o importato, entra in gioco soprattutto il servizio asset; quando invece parte un turno conversazionale, l'audio viene trascritto da Whisper, il testo passa dal servizio RAG e la risposta viene sintetizzata dal TTS con la voce associata. Distinguere questi due percorsi aiuta a capire perché alcune criticità riguardino la latenza del turno e altre la consistenza del profilo tra sessioni diverse.

Lo stesso vale per la ricerca ibrida usata dal RAG: i candidati vengono ordinati combinando similarità vettoriale e BM25, con un leggero favore per la componente lessicale. La scelta nasce dal tipo di query che il prototipo riceve, spesso ricche di nomi propri, parole chiave e riferimenti espliciti, mentre il recupero vettoriale resta utile quando la memoria viene richiamata in forma parafrasata. La Figura~\ref{fig:rag-swagger-cap3} riporta una vista della documentazione automatica degli endpoint: serve a visualizzare la granularità delle route, non a sostituire la descrizione architetturale proposta qui. Una ricostruzione più analitica dei flussi, dei dati scambiati e delle responsabilità dei servizi è raccolta nelle Appendici~\ref{app:rag} e \ref{app:backend-services}.

\begin{figure}[!htbp]
\centering
\includegraphics[width=1.0\textwidth]{rag_docs_swagger.png}
\caption{Interfaccia Swagger UI (\texttt{/docs}) del servizio RAG di
SOULFRAME: una vista della documentazione automatica generata da FastAPI per \texttt{rag\_server.py}.}
\label{fig:rag-swagger-cap3}
\end{figure}

\subsection{Persistenza e gestione dati per avatar}

Nel sistema, la persistenza dei dati di profilo è organizzata su filesystem e separata per \texttt{avatar\_id}: ogni avatar mantiene tre store dedicati, \texttt{avatar\_store/} per gli asset \texttt{.glb}, \texttt{voices/avatars/} per il profilo vocale e \texttt{rag\_store/} per la memoria conversazionale persistente. In questo modo le risorse operative restano allineate lungo tutto il ciclo di vita del profilo. La Figura~\ref{fig:persistenza-tree-cap3} mostra la distribuzione dei tre store persistenti e i punti di verifica disponibili.

\begin{figure}[!htbp]
\centering
\includegraphics[width=0.95\textwidth]{persistenza_tree_cap3.png}
\caption{Struttura del filesystem di persistenza in SOULFRAME:
tre store isolati per \texttt{avatar\_id} (asset \texttt{.glb},
profilo vocale e memoria RAG).}
\label{fig:persistenza-tree-cap3}
\end{figure}

L'isolamento per profilo consente di evitare interferenze tra avatar distinti e di ricostruire lo stato dopo riavvio rileggendo i rispettivi percorsi persistenti. La figura distingue questi store stabili dai dati di sessione: nel servizio RAG i turni recenti vengono mantenuti in RAM per coppia \texttt{(avatar\_id, session\_id)}, mentre i log conversazionali sono scritti separatamente su file quando il client apre una sessione con \texttt{/chat\_session/start}. Più precisamente, \texttt{rag\_store} conserva la memoria riutilizzabile dell'avatar, il contesto recente serve al turno corrente e il log documenta la sessione. Nell'applicazione, la presenza di memoria può essere verificata tramite \texttt{/avatar\_stats}, che segnala se l'avatar ha già una memoria popolata, mentre asset e voce restano associati allo stesso \texttt{avatar\_id} nei rispettivi store.

Dal punto di vista del lavoro sperimentale, questa separazione ha avuto un vantaggio molto pratico: permettere controlli più chiari sui guasti. Se manca memoria, il problema resta nel ramo RAG; se manca la voce, il controllo cade sul TTS; se un modello non viene più caricato, il problema resta nel catalogo asset. In una pipeline del genere, sapere dove inizia e dove finisce ogni persistenza conta quasi quanto la persistenza stessa.


\FloatBarrier

\section{Setup e deploy operativo}
Il setup e il deploy operativo traducono l'architettura descritta nella Sezione~\ref{sec:architettura-riferimento} in procedure ripetibili di avvio, arresto e verifica. SOULFRAME mantiene la stessa organizzazione logica dal contesto di sviluppo a quello pubblico, ma cambia il modo in cui i processi vengono orchestrati e raggiunti in rete. Per ridurre gli errori di esercizio, sono stati privilegiati passaggi ripetibili e un controllo unificato dei servizi.

\subsection{Ambienti di esecuzione e deploy}
In ambiente locale la configurazione è orientata a sviluppo e test rapidi: i servizi girano sull'indirizzo della macchina e il frontend WebGL è servito in HTTP. Il provisioning iniziale è supportato da \texttt{setup\_soulframe\_windows.bat}, mentre l'avvio operativo è centralizzato in \texttt{ai\_services.cmd}, così da mantenere un punto unico per start, stop e restart dello stack locale. Questo assetto favorisce diagnostica iterativa e osservabilità immediata; le differenze rispetto al deploy server sono sintetizzate in Tabella~\ref{tab:deploy-confronto-cap3}.

Nel lavoro pratico questa centralizzazione ha avuto un vantaggio molto concreto: ridurre il tempo tra modifica del codice e verifica end-to-end. Restano comunque esterni alcuni prerequisiti di sistema, in particolare Ollama e Tesseract, che in locale incidono direttamente sulla ripetibilità del setup.

\begin{table}[!htbp]
\centering
\footnotesize
\setlength{\tabcolsep}{3pt}
\renewcommand{\arraystretch}{1.2}
\begin{tabular}{|>{\raggedright\arraybackslash}p{0.19\textwidth}|>{\raggedright\arraybackslash}p{0.35\textwidth}|>{\raggedright\arraybackslash}p{0.35\textwidth}|}
\hline
\textbf{Aspetto} & \textbf{Windows locale} & \textbf{Ubuntu server} \\
\hline
Avvio servizi & Script locale (\texttt{ai\_services.cmd}) con processi separati in ascolto su loopback & Service manager con unità dedicate e restart automatico \\
\hline
Origine/URL usati dal client & Origine locale (\texttt{http://localhost:8000}) con porte backend esplicite in configurazione client & Unica origine HTTPS; API raggiunte tramite percorsi \texttt{/api/<servizio>} dietro reverse proxy \\
\hline
TLS & Assente nel loop di sviluppo locale (HTTP) & Terminazione TLS al reverse proxy \\
\hline
Logging & Log di console/processo per debugging interattivo & Log centralizzati di servizi e reverse proxy \\
\hline
Update/rollback & Aggiornamento manuale dello workspace e riavvio script & Script amministrativi con stop/start orchestrato e supporto backup/rollback \\
\hline
Gestione persistenze & Directory locali del progetto (\texttt{avatar\_store}, \texttt{voices}, \texttt{rag\_store}) & Directory persistenti lato server con separazione per componente/avatar \\
\hline
\end{tabular}
\caption{Differenze operative essenziali tra setup locale Windows e deploy server Ubuntu.}
\label{tab:deploy-confronto-cap3}
\end{table}

\paragraph{Deploy su Ubuntu e servizi di supporto.}
In ambiente Ubuntu l'obiettivo è l'esposizione pubblica controllata del prototipo. I micro-servizi restano su rete interna, mentre il browser raggiunge un'unica origine HTTPS tramite reverse proxy. Il setup server separa componenti applicativi, file statici WebGL e variabili ambientali operative; la gestione dei processi è demandata a unità del service manager con politiche di restart, così che il ripristino dopo reboot o failure non richieda interventi manuali su ogni servizio.

Accanto ai micro-servizi applicativi, il prototipo include un livello di supporto operativo per gestione dei processi, reverse proxy e automazione amministrativa. In particolare, i servizi restano in bind su \texttt{127.0.0.1} e il routing pubblico viene concentrato in un solo punto, mentre helper come \texttt{sfctl} e \texttt{sfurl} rendono più semplice distinguere un problema applicativo da un problema di deploy o di instradamento.

La complessità operativa non scompare, ma diventa più leggibile. Nel passaggio dal locale al server non cambia solo il sistema operativo: cambia anche il modo in cui lo stack viene avviato, osservato e mantenuto nel tempo. In una tesi centrata su un prototipo funzionante, il deploy non è quindi un dettaglio secondario ma una parte concreta della tenuta complessiva del progetto.

Il frontend Unity WebGL resta così riutilizzabile nei due contesti, variando soprattutto il livello di orchestrazione e il punto di ingresso di rete. La tabella comparativa precedente rende verificabili queste differenze; i dettagli esecutivi sono approfonditi nel Capitolo~\ref{chap:implementazione} e in Appendice~\ref{app:ops}.
